\def\level{Première}  % Classe à droite
\def\course{NSI}
\def\chaptername{TP Réseau Packet Tracer} % Titre au centre
\def\docname{ARSE05-TP1}        % Identifiant à gauche

\pagestyle{cours_FP}
\makeheader{} % Affiche le bandeau


\begin{methode}{}{}
\begin{itemize}
    \item Lancer le logiciel \verb+Packet Tracer+ (Bureau / Mathématique en Local)
    \item Ajouter un poste dans le réseau: \verb+Terminaux / Generic+. Le poste ajouté s'appelle \verb+PC-PT PC0+
    \item pour supprimer ce poste ou tout autre élément vous pouvez cliquer sur l'icône \verb+Delete(Del)+ à droite de l'écran puis sélectionner l'élément à supprimer.
\end{itemize}
\end{methode}

\section{Communication poste à poste}

\begin{exercice_cours}[Connexion poste à poste]{}

    \medskip

    \begin{minipage}{0.60\linewidth}
\begin{enumerate}
            \item Placer deux postes \verb+Generic+ nommés PC01 et PC02
            \item Relier les deux postes par une connexion automatique. 
\includegraphics[width = 0.05\linewidth]{./icone1.jpg}
\end{enumerate}       
    \end{minipage}\hfill
    \begin{minipage}{0.38\linewidth}
        \begin{center}
                    \includegraphics[width = 0.8\linewidth]{./icone02.jpg}
        \end{center}
    \end{minipage}
\begin{remarque}{}{}
En choisissant la connexion automatique, le logiciel se charge de choisir le meilleur type de connexion. Ici la liaison utilisée est un câble croisé.
\end{remarque}

\end{exercice_cours}

\begin{exercice_cours}[Communication entre les postes]{}

    Pour faire communiquer ces deux postes on peut utiliser un \verb+PDU+ (Protcol Data Unit).

    \begin{enumerate}
        \item Cliquer sur l'icône ``Add simple PDU'' 
\includegraphics[width = 0.05\linewidth]{./icone03.jpg}
        \item Cliquer dans l'ordre sur le poste émetteur puis sur le destinataire.
        \item Vous obtenez un message d'erreur: \verb+PCI has no functional ports+
        \item Expliquer pourquoi les informations ne peuvent pas circuler.
        
        \begin{tcolorbox}[height = 1cm]
        
        \end{tcolorbox}

        \begin{minipage}{0.75\linewidth}
        \item Nous allons définir les adresses IP pour chaque poste. Cliquer sur l'icône ``Select'' puis sur l'un des postes.
        
        Choisir ``Desktop / IP configuration'' et taper l'adresse \verb+192.168.1.1+. Le masque de sous réseau va se définir automatiquement (classe C). 
    
        \item Faire de même avec l'autre poste et lui attribuer l'adresse \verb+192.168.1.2+
        
        \item Recommencer l'envoi d'une trame entre ces deux poste. Cette fois-ci la transmission doit se dérouler normalement.
         
        \end{minipage}\hfill
        \begin{minipage}{0.23\linewidth}
            \includegraphics[width = 1\linewidth]{./icone04.jpg}
        \end{minipage}   
        
        \medskip

        La communication s'est déroulée très vite et nous n'avons pas eu le temps de voir ce qu'il se passait. 

        \item Passer en mode \verb+Simulation+ en cliquant sur l'icône en bas à droite. 
    
        \item Cliquer sur ``Edit Filters'' puis décocher l'option ``show all'' et cocher uniquement le protocole ICMP. Nous ferons cette manipulation à chaque fois que nous utiliserons ce mode dans ce TP.
    
        \item cliquer sur ``Auto Capture/Play'' et observer l'animation entre les deux postes. Pour réinitialiser la simulation cliquer en bas sur ``Delete''
        
        \item Décrivez ce que vous voyez
        
        \begin{tcolorbox}[height = 1cm]
        
        \end{tcolorbox}

    \end{enumerate} 

\end{exercice_cours}

\section{Faire communiquer plusieurs réseaux entre eux}

\begin{exercice_cours}[A l'aide du réseau filaire]{}



\begin{minipage}{0.65\linewidth}
    Nous allons simuler un réseau composé de 4 ordinateurs et d'un switch (commutateur). Réaliser le schéma suivant puis affecter l'adressage IP:

    \begin{itemize}
        \item \verb+PC0: 192.168.1.1/24+
        \item \verb+PC1: 192.168.1.2/24+
        \item \verb+PC2: 192.168.1.3/24+
    \end{itemize}
\end{minipage}\hfill
\begin{minipage}{0.33\linewidth}
       \includegraphics[width = 0.9\linewidth]{./icone07.jpg}
 
\end{minipage}

    \begin{enumerate}
        \item Envoyer un message (add Simple PDU) de PC0 vers PC2 en respectant le consignes du précédent exercice (Simulation / ICMP / Capture automatique/ Jouer). Décrire ce qu'il se passe.
        
        \begin{tcolorbox}[height = 2cm]
        
        \end{tcolorbox}

        \item Pour tester la communication entre plusieurs postes au sein d'un réseau, il est possible d'utiliser la commande \verb+ping+ (Packet INternet Groper). Cette commande permet de tester l'accessibilité à une autre machine (adresse IP). \verb+ping+ utilise une requête ICMP \verb+echo+ et attend en retour une réponse \verb+echo reply+.
        
        Tester la communication des 4 PC entre eux avec la commande. ``Desktop / command prompt''

        Dans la fenêtre qui apparait après le prompt \verb+C:\+, taper la commande \verb+ping <adresse IP>+ avec l'adresse que vous voulez joindre.

        Faites un test de communication pour chaque PC. Qu'observe-t-on? Indiquer ce que fait chaque requête \verb+echo+

        \begin{tcolorbox}[height = 3cm]
        
        \end{tcolorbox}

        \item Modifier l'adresse IP de PC2 en \verb+192.168.6.4/24+ et refaite un \verb+ping+ depuis PC0 vers PC3. Que constatez-vous? Expliquer.
       
        \begin{tcolorbox}[height = 3cm]
        
        \end{tcolorbox}

        \item Redonner à PC2 l'adresse \verb+192.168.1.3+. Relever les adresses \verb+MAC+ des 3 cartes Ethernet des PC:
        \begin{itemize}
            \item PC0: \dotfill
            \item PC1: \dotfill
            \item PC2: \dotfill
        \end{itemize}

        \item Tester la communication entre les PC de votre salle à l'aide de la commande \verb+ping+. Tester également d'autres adresses IP de PC d'autres salles.
        
    \end{enumerate}

\end{exercice_cours}

\medskip

\begin{exercice_cours}[Communication à l'aide d'un réseau sans fil]{}

\medskip

        \begin{minipage}{0.7\linewidth}
        Nous allons simuler un réseau sans fil WiFi (Wireless Fidelity) constitué de:
            \begin{itemize}
                \item un ordinateur portable (Laptop-PT Generic)
                \item un serveur (server-PT generic)
                \item un point d'accès WiFi (access point-PT Generic)
            \end{itemize}

            Aucun de ces terminaux ne peut se connecter au point d'accès car ils n'ont pas de connexion WiFi
        \end{minipage}\hfill
        \begin{minipage}{0.28\linewidth}
            \includegraphics[width = 0.9\linewidth]{./icone08.jpg}

        \end{minipage}
        
        \medskip

        \begin{enumerate}
   
            \item Cliquer sur le serveur puis onglet \verb+Physique+ sur la face arrière visible, éteindre le serveur et remplacer le connecteur Fast Ethernet par un connecteur Wireless (PT-HOST-NM-1W). Rallumer le serveur.

            \item Donner au serveur une IP fixe: ``Configuration, Wireless0'' et dans la rubrique ``IP configuration'', cocher ``Static'' puis entrer l'adresse IP \verb+192.168.1.253/24+. Fermer la fenêtre.
            
            \item Configurer l'ordinateur portable en remplaçant le connecteur. Mettre sa configuration IP en mode DHCP.
            \item Rajouter deux autres périphériques sans fils: une tablette PC et un PDA.
    
            \item Observer les adresses IP des différents périphériques. Quelle remarque pouvez vous faire?
        
            \begin{tcolorbox}[height = 2cm]
        
            \end{tcolorbox}
    
        \end{enumerate}

\end{exercice_cours}

\medskip

\begin{exercice_cours}[Communication entre deux réseaux]{}

    \medskip

    \begin{enumerate}

        \begin{minipage}{0.55\linewidth}
            \item Réaliser le réseau ci-contre de manière à avoir deux sous réseaux avec une adresse IP différente:

        \begin{itemize}
            \item A gauche: sous réseau d'adresse \verb+192.168.1.0/24+
            \item A droite: sous réseau d'adresse \verb+172.16.0.0/16+
        \end{itemize}

        N'oubliez pas d'attribuer les adresses IP fixes.
    
        \end{minipage}\hfill
        \begin{minipage}{0.43\linewidth}
            \includegraphics[width = 0.9\linewidth]{./icone09.jpg}
 
        \end{minipage}
        
        \item Connecter les deux switchs entre eux avec un câble et tester la connexion entre deux PC de réseaux différents. Que constatez-vous? Expliquer.
        
        \begin{tcolorbox}[height = 2cm]
        
        \end{tcolorbox}

        \item Supprimer ce câble et insérer un routeur entre les deux switchs (routeur 1841). Reliez les sous réseaux comme sur le schéma. Testez les communications entre PC0 et PC3. Expliquer.
        
        \begin{minipage}{0.48\linewidth}
           \begin{center}
                    \includegraphics[width = 0.95\linewidth]{./icone10.jpg}

        \end{center} 
        \end{minipage}\hfill{}
                \begin{minipage}{0.48\linewidth}
\begin{tcolorbox}[height = 3cm]
        
        \end{tcolorbox}
        \end{minipage}
        


        

        \medskip

        \begin{easybox}[Configuration des adresses IP du routeur]{}
        La configuration du routeur doit respecter les contraintes suivantes:
        \begin{itemize}
            \item le port du routeur relié au sous réseau de gauche doit posséder une adresse IP de ce sous réseau. En général on lui donne la dernière adresse disponible.
            \item même configuration pour le port du routeur relié au sous réseau de droite.
        \end{itemize}
        \end{easybox}

        \item Attribuez des adresses IP aux ports du routeur. Indiquer sur le schéma les différentes adresses IP que vous avez choisi. 
        \item Tester maintenant les connexions entre PC0 et PC3. Que remarquez-vous? Expliquer
        
        \begin{tcolorbox}[height = 1.5cm]
        
        \end{tcolorbox}

        \begin{easybox}[Configurer la passerelle par défaut]{}
            Le réseau de gauche et le réseau de droite ne peuvent pas encore communiquer car les paquets ne savent pas par où sortir du sous réseau.

            Pour chaque machine, la passerelle par défaut (Default Gateway) doit être configurée pour pointer vers l'adresse IP du port du routeur correspondant à son sous-réseau. Par exemple, pour le sous-réseau de gauche (\verb+192.168.1.0/24+), la passerelle par défaut sera l'adresse IP du port du routeur connecté à ce sous-réseau. De même, pour le sous-réseau de droite (\verb+172.16.0.0/16+), la passerelle par défaut sera l'adresse IP du port du routeur connecté à ce sous-réseau.
        \end{easybox}

        \item Configurer les passerelles par défaut des chaque PC.
        
        \item Tester maintenant la communication entre PC0 et PC3. Que constatez-vous?
        
        
        \begin{tcolorbox}[height = 1.5cm]
        
        \end{tcolorbox}
    \end{enumerate}
\end{exercice_cours}

\medskip

\begin{exercice_cours}[Serveur HTTP et DNS]{}

    \medskip

    Nous allons simuler deux choses:
    \begin{itemize}
        \item l'hébergement d'un site internet sur un serveur HTTP
        \item la redirection des adresses IP grâce au serveur DNS
    \end{itemize}

    \begin{enumerate}
        \item Créer le réseau suivant en configurant comme sur le schéma les différents objets connectés.
        
        \begin{center}
        \includegraphics[width = 0.75\linewidth]{./icone12.jpg}
            
        \end{center}

        \item Donner les valeurs suivantes dans le tableau:
        
        \begin{center}
        \begin{tabular}{|c|c|c|c|}
        \hline
        Nom du périphérique & Adresse IP & Masque de sous réseau & Passerelle \\
        \hline
        PC\_A &  &  &  \\ \hline
        Routeur coté réseau A & & & XXX \\ \hline
        PC\_B &  &  &  \\ \hline
        Routeur coté réseau B & & & XXX\\ \hline
        Serveur HTTP &  &  &  \\ \hline
        Serveur DNS &  &  &  \\ \hline
        \end{tabular}
        \end{center}

        \item Quelle est l'adresse du réseau A, son adresse de broadcast? \dotfill{}
        \item Quelle est l'adresse du réseau B, son adresse de broadcast? \dotfill{}
        \item Vérifier la communication entre les différents éléments du réseau. A ce stade il ne devrait pas y avoir de problème.
        \item Cliquer sur le serveur HTTP puis ``HTTP''. Vérifier que le protocole HTTP est bien ``On''.
        \medskip

        Trois sites sont présents sur le serveur. Nous travaillerons avec \verb+http://www.index.html+ mais vous pouvez choisir un des deux autres.
        \item Vérifier que le site web est accessible en ouvrant ``PC\_A / Desktop / Web Browser'' et en tapant l'adresse IP du serveur HTTP.
        \item Faites pareil avec l'adresse \verb+http://www.index.html+. Que constatez vous? Expliquer.
        
        \begin{tcolorbox}[height = 2cm]
        
        \end{tcolorbox}

        \item Pour faire la ``traduction'' entre les adresses IP et les noms de domaine, il faut renseigner l'adresse du serveur DNS pour les PC\_A et PC\_B ainsi que pour le serveur HTTP.
        
        Cliquer sur chaque périphérique puis ``Ip configuration / DNS Server'' et renseigner l'adresse IP du serveur DNS
    
        \item Tester maintenant la connexion au site \verb+http://www.index.html+. Que constatez vous? Expliquer.
        
        \begin{tcolorbox}[height = 2cm]
        
        \end{tcolorbox}

        \item Il faut faire le lien entre l'adresse IP du serveur et le nom de domaine
        
        Dans le serveur DNS, cliquer sur ``Config / DNS'' et renseigner le nom du site: \verb+http://www.index.html+ et l'adresse du serveur HTTP. Cliquer sur ``Add''

        \begin{center}
        \includegraphics[width = 0.5\linewidth]{./icone13.jpg}
        \end{center}

        \item Tester à nouveau la connexion au site \verb+http://www.index.html+
        
        \item Dans le serveur DNS rajouter une ligne en tapant \verb+www.banquepopulaire.fr+ comme nom et l'adresse du serveur HTTP. Cliquer sur ``Add''
        

        \item Tester alors la connexion au site \verb+www.banquepopulaire.fr+
        
        \begin{tcolorbox}
        Vous venez d'expérimenter de manière simple le principe de l'attaque DNS. En d'autres termes, si quelqu'un parvient à manipuler les enregistrements DNS pour rediriger l'utilisateur vers des sites frauduleux, celui ci pourrait vous amener à entrer des données sensibles.
        \end{tcolorbox}
    \end{enumerate}
\end{exercice_cours}


    \medskip

    \begin{exercice_cours}[Mettre en place un serveur ftp]{}

        \begin{definition}{Le serveur ftp}{}
            Le serveur \verb+ftp+ (File Transfert Protocol) permet de transférer des fichiers par Internet. Toute personne ayant l'autorisation peut télécharger ou envoyer des fichiers depuis un ordinateur distant.
        \end{definition}

        \begin{enumerate}
            \item Dans le serveur HTTP, cliquer sur ``Config / FTP''. Un utilisateur est déjà renseigné. 
            
            Noter l'identifiant et le mot de passe: \dotfill
            \item Sur un des PC, ouvrir l'invite de commande et connectez vous au serveur \verb+ftp+ en tapant \verb+ftp 172.16.0.15+. Entrer l'identifiant et le mot de passe.
            \item La commande \verb+dir+ permet de lister les fichiers sur le serveur \verb+ftp+
            \item Pour ajouter un fichier vous pouvez utiliser la commande \verb+put+. A la racine de votre machine virtuelle vous avez un fichier nommé ``\verb+sampleFile.txt+''. Transférez le sur le serveur \verb+ftp+ avec ``\verb+put sampleFile.txt+''.
            \item Vérifier que le transfert s'est bien passé en listant les fichiers présents avec \verb+dir+.
            \item Récupérez un autre fichier depuis le serveur \verb+ftp+ avec ``\verb+get <nom du fichier>+''.
            \item Quitter le serveur \verb+ftp+ en faisant \verb-Control + C-.
            \item Vérifier que vous avez bien récupéré le fichier sur le disque local de la machine avec la commande \verb+dir+.
        \end{enumerate}
    \end{exercice_cours}

